\section{Numerical Methods}
This chapter contains an overview of the governing equations used in Star-CCM+ for the modelling of multiphase flows with a focus on
eulerian multiphase (EMP) flow and the volume of fluid (VOF) method. Turbulence and the Shear Stress Transport (SST) k-$\omega$ model is
briefly discussed. The physical processes of evaporation and vaporisation are described.
The ability to model multiple flow regimes, capture interfaces, and the phase interaction closures used in Star-CCM+ are discussed.
The chapter concludes with a description of the numerical methods used in Star-CCM+ along with the apparent limitations~\COMMENT{better word}
of the commercial software in being able to accurately model the problems under consideration.

\subsection{Governing Equations}
\COMMENT{define region as in what is V and what is delV and stuff}
Multiphase flow is modelled in Star-CCM+ using various frameworks. The two-fluid - eulerian multiphase - framework solves a set of six equations
to model two-phase flow. A set of transport equations for mass, momentum, and energy are solved for all the phases, which are assumed to
be interpenetrating continua that share a common pressure field. The local instantaneous conservation equations are assumed to be euler averaged and
contain closures based on interfacial terms that define the influences of the phases upon each other.

Although used primarily for modelling dispersed flows, the EMP framework is capable
of resolving both dispersed and separated flow regimes.

The implemented volume of fluid method in Star-CCM+ assumes sufficient mesh resolution
in order to resolve the interface and model separated flows. The interface capturing method solves
the volume fraction transport equation using the High-Resolution Interface Capturing (HRIC) scheme
to maintain a sharp interface between the phases.

Based on the volume fraction of the phase $i$, $\alpha_{i} = V_{i}/V$, where in each cell $\textstyle\sum\nolimits_{i} \alpha_{i} = 1$.
The material properties in each cell and other flow variables are assumed to be a mixture of the properties of the phases as follows:
\begin{subequations}\label{eq_vof_defs}
\begin{align}
\text{density, } \rho &= \textstyle\sum\nolimits_{i} \alpha_{i} \rho_{i},\\
\text{viscosity, } \mu &= \textstyle\sum\nolimits_{i} \alpha_{i} \mu_{i},\\
\text{velocity, } \vetor{v} &= \textstyle\sum\nolimits_{i} \alpha_{i} \rho_{i} \vetor{v}_{i}/ \rho,\\
\text{specific heat, } C_{p} &= \textstyle\sum\nolimits_{i} \alpha_{i} C_{p,i} \rho_{i}/\rho,\\
\text{thermal conductivity, } k &= \textstyle\sum\nolimits_{i} \alpha_{i} k_{i},\\
\text{enthalpy, } \mathrm{H} &= \textstyle\sum\nolimits_{i} \alpha_{i} \rho_{i} \mathrm{H}_{i}/\rho + |\vetor{v}|^{2}/2.
\end{align}
\end{subequations}

\noindent where $\rho_{i}$ is the density, $\mu_{i}$ is the dynamic viscosity,
$\vetor{v}_{i}$ is the velocity, $C_{p,i}$ is the specific heat, $k$ is the thermal conductivity, and $\mathrm{H}_{i}$ is the enthalpy of the phase $i$.

The diffusion velocity of phase $i$, $\vetor{v}_{d,i}$, is defined as the difference between the velocity of the phase and the mixture
such that \[\vetor{v}_{d,i} = \vetor{v}_{i} - \vetor{v}.\]

A set of volume fraction transport and mass, momentum, and energy conservation equations are solved for each phase.
In cases of two-phase flow, Star-CCM+ solves the volume fraction transport equation only for the first phase.

The governing equations used in the implementation of both the EMP and VOF frameworks in
Star-CCM+ are summarised below.

\subsubsection{Eulerian Multiphase Flow}

The \textbf{mass conservation} equation solved for the phase $i$ is

\begin{equation}\label{eq_emp_cont}
\frac{\partial}{\partial t}\int\limits_{V} \alpha_{i} \rho_{i}\,\mathrm{d}V +
\int\limits_{\partial V} \alpha_{i} \rho_{i} \vetor{v}_{i} \cdot \mathrm{d}\vetor{a}
= \int\limits_{V} \sum_{j \neq i} (m_{ij} - m_{ji})\,\mathrm{d}V + \int\limits_{V} \mathrm{S}^{\alpha}_{i}\,\mathrm{d}V
\end{equation}
\COMMENT{CHECK}%\hspace*{5pt} 
where for the phase $i$
\begin{itemize}[topsep=0pt, labelindent=40pt,labelsep=5pt, parsep = 2pt, listparindent=2em]
\item $\vetor{a}$ is the surface area vector,
\item $\alpha_{i}$ is the volume fraction that fulfills the requirement $\textstyle\sum\nolimits_{i}^{n} \alpha_{i} = 1$,
\item $\rho_{i}$ is the density,
\item $\vetor{v}_{i}$ is the velocity,
\item $m_{ij}, m_{ji}$ is the interphase mass transfer per unit volume from phase $i$ to $j$ and $j$ to $i$ respectively,
\item $\mathrm{S}^{\alpha}_{i}$ is the user defined mass source term.
\end{itemize}


The \textbf{momentum conservation} equation solved for the phase $i$ is
\begin{equation}\label{eq_emp_mom}
\begin{split}
\frac{\partial}{\partial t}\int\limits_{V} \alpha_{i} \rho_{i} \vetor{v}_{i}\,\mathrm{d}V +
\int\limits_{\partial V} \alpha_{i} \rho_{i} \vetor{v}_{i} \otimes \vetor{v}_{i} \cdot \mathrm{d}\vetor{a}
&= -\int\limits_{V} \alpha_{i} \nabla p\,\mathrm{d}V + \int\limits_{V} \alpha_{i} \rho_{i} \vetor{g}\,\mathrm{d}V
+ \int\limits_{\partial V}\left[\alpha_{i} \left(\tens{T}_{i} + \tens{T}^{t}_{i}\right)\right]\cdot \mathrm{d}\vetor{a}
+ \int\limits_{V} \vetor{M}_{i}\,\mathrm{d}V \\
&\quad + \int\limits_{V}\left(\vetor{F}_{\text{int}}\right)_{i}\,\mathrm{d}V
+ \int\limits_{V}\vetor{S}^{\alpha}_{i}\,\mathrm{d}V 
+ \int\limits_{V} \sum_{j \neq i} (m_{ij}\vetor{v}_{j} - m_{ji}\vetor{v}_{i})\,\mathrm{d}V
\end{split}
\end{equation}

where for the phase $i$
\begin{itemize}[topsep=0pt, labelindent=40pt,labelsep=5pt, parsep = 2pt, listparindent=2em]
\item $p$ is the pressure,
\item $\vetor{g}$ is the gravitational force vector,
\item $\tens{T}_{i}, \tens{T}^{t}_{i}$ are the molecular and turbulent stresses respectively,
\item $\vetor{M}_{i}$ is the interphase momentum transfer term per unit volume that satisfies $\textstyle\sum\nolimits_{i = 1}^{n}\vetor{M}_{i} = 0$,
\item $\left(\vetor{F}_{\text{int}}\right)_{i}$ is the internal force vector that consists of either the
solid pressure force between particles or user defined potential force,
\item $\vetor{S}^{\alpha}_{i}$ is the user defined momentum source term.
\end{itemize}


The \textbf{energy conservation} equation solved for the phase $i$ is
\begin{equation}\label{eq_emp_energy}
\begin{split}
\frac{\partial}{\partial t}\int\limits_{V} \alpha_{i} \rho_{i} \mathrm{E}_{i}\,\mathrm{d}V +
\int\limits_{\partial V} \alpha_{i} \rho_{i} \mathrm{H}_{i} \vetor{v}_{i} \cdot \mathrm{d}\vetor{a} &=
\int\limits_{\partial V} \alpha_{i} (k_{\mathrm{eff}})_{i} \nabla T_{i} \cdot \mathrm{d}\vetor{a} + \int\limits_{\partial V} \tens{T}_{i} \vetor{v}_{i} \cdot \mathrm{d}\vetor{a}
+ \int\limits_{V} \vetor{f}_{i} \cdot \vetor{v}_{i}\,\mathrm{d}V +\int\limits_{V}\mathrm{S}_{u,i}\,\mathrm{d}V \\
&{} + \int\limits_{V} \sum_{j \neq i} Q_{ij}\,\mathrm{d}V +
\int\limits_{V} \sum_{ij} Q_{i}^{(ij)}\,\mathrm{d}V 
+ \int\limits_{V} \sum_{j \neq i} (m_{ij} - m_{ji})h_{i}(T_{ij})\,\mathrm{d}V
\end{split}
\end{equation}

\noindent where for the phase $i$,
\begin{itemize}[topsep=0pt, labelindent=40pt,labelsep=5pt, parsep = 2pt, listparindent=2em]
\item $\mathrm{E}_{i}$ is the total energy,
\item $\mathrm{H}_{i}$ is the total enthalpy,
\item $\tens{T}_{i}$ is the viscous stress tensor,
\item $T_{i}$ is the temperature,
\item $\vetor{f}_{i}$ is the body force vector,
\item $(k_{\mathrm{eff}})_{i}$ is the effective thermal conductivity,
\item $\mathrm{S}_{u,i}$ is the user defined energy source term,
\item $Q_{ij}$ is the heat transfer from phase $i$ to phase $j$ due to temperature difference between the phases,
\item $Q_{i}^{(ij)}$ is the heat transfer between the phase $i$ and the interface which accounts for phenomenon
such as evaporation and boiling,
\item $h_{i}(T_{ij})$ is the enthalpy of the phase $i$ evaluated at the temperature $T_{ij}$ of the interface.
\end{itemize}

The effective thermal conductivity of the phase $i$ is defined using the thermal conductivity, $k_{i}$,
the turbulent viscosity, $\mu_{t,i}$, the specific heat, $C_{p,i}$, and the turbulent thermal diffusion Prandtl number,
$\sigma_{t,i}$ as \[(k_{\mathrm{eff}})_{i} = k_{i} + \frac{\mu_{t,i} C_{p,i}}{\sigma_{t,i}}\].

The total enthalpy of the phase $i$, $\mathrm{H}_{i}$, is defined using the specific enthalpy $h_{i}$
and the kinetic energy as \[\mathrm{H}_{i} = h_{i} + \frac{1}{2}\vetor{v}_{i}\cdot\vetor{v}_{i},\]

where the specific enthalpy is further defined using the specific heat of the phase and the reference temperatures
and heat of formation such that

\[h_{i} = h_{i}^{\text{REF}} + \int\limits_{T_{i}^{\text{REF}}}^{T_{i}} C_{p,i} (T')\,\mathrm{d}T'.\]

The total energy of the phase $i$ is then related to the total enthalpy as $\mathrm{E}_{i} = \mathrm{H}_{i} - p/\rho_{i}$.

The interphase energy terms are related as follows:

\begin{itemize}
\item for cases where heat is exchanged only through thermal diffusion, the interphase heat transfer rates satisfy $Q_{ij} = -Q_{ji}$
\item for cases where heat and mass are transferred only through phase change, the interphase heat and mass transfer rates satisfy
$Q_{i}^{(ij)} + Q_{j}^{(ij)} + \left(m_{ij} - m_{ji}\right)\left(h_{j}(T_{ij}) - h_{i}(T_{ij})\right) = 0$.
\end{itemize}
\subsubsection{Volume of Fluid Method}
The distribution of the phase $i$ is driven by the volume fraction transport.
\begin{equation}\label{eq_vof_transport}
\frac{\partial}{\partial t}\int\limits_{V} \alpha_{i}\,\mathrm{d}V +
\int\limits_{\partial V} \alpha_{i} \vetor{v}_{i} \cdot \mathrm{d}\vetor{a}
= \int\limits_{V}\left(\mathrm{S}^{\alpha}_{i} - \frac{\alpha_{i}}{\rho_{i}} \frac{D \rho_{i}}{D t}\right)\,\mathrm{d}V
- \int\limits_{V} \frac{1}{\rho_{i}}\nabla (\alpha_{i}\rho_{i}\vetor{v}_{d,i})\,\mathrm{d}V
\end{equation}
\noindent where $\vetor{a}$ is the surface area vector, $\vetor{v}$ is the mixture velocity, $\vetor{v}_{d,i}$ is the slip velocity, and
$\mathrm{S}^{\alpha}_{i}$ is the user defined phase source term for the phase $i$.


The \textbf{mass conservation} equation solved is
\begin{equation}\label{eq_vof_mass}
\frac{\partial}{\partial t}\int\limits_{V} \rho\,\mathrm{d}V +
\int\limits_{\partial V} \rho \vetor{v} \cdot \mathrm{d}\vetor{a}
= \int\limits_{V} \sum_{i}\mathrm{S}^{\alpha}_{i} \rho_{i}\,\mathrm{d}V
\end{equation}


The \textbf{momentum conservation} equation solved is
\begin{equation}\label{eq_vof_mom}
\begin{split}
\frac{\partial}{\partial t}\int\limits_{V} \rho \vetor{v}\,\mathrm{d}V +
\int\limits_{\partial V} \rho\vetor{v} \otimes \vetor{v} \cdot \mathrm{d}\vetor{a}
&= -\int\limits_{V} \nabla p\,\mathrm{d}V + \int\limits_{V} \rho \vetor{g}\,\mathrm{d}V
+ \int\limits_{\partial V}\left[\left(\tens{T} + \tens{T}^{T}\right)\right]\cdot \mathrm{d}\vetor{a} 
 + \int\limits_{V} \vetor{f}_{b}\,\mathrm{d}V\\
& - \sum_{i}\int\limits_{\partial V} \alpha_{i} \rho_{i} \vetor{v}_{d,i} \otimes \vetor{v}_{d,i} \cdot \mathrm{d}\vetor{a}
 + \int\limits_{V}\vetor{S}^{\alpha}_{i}\,\mathrm{d}V
\end{split}
\end{equation}

\noindent where
\begin{itemize}[topsep=0pt, labelindent=40pt,labelsep=5pt, parsep = 2pt, listparindent=2em]
\item $p$ is the pressure,
\item $\vetor{g}$ is the gravitational force vector,
\item $\tens{T}$ is the viscous stress tensor,
\item $\vetor{f}_{b}$ is the body force vector,
\item $\vetor{S}^{\alpha}_{i}$ is the user defined phase momentum source term.
\end{itemize}


The \textbf{energy conservation} equation solved is
\begin{equation}\label{eq_vof_energy}
\begin{split}
\frac{\partial}{\partial t}\int\limits_{V}\rho \mathrm{E}\,\mathrm{d}V +
\int\limits_{\partial V} \left[\rho \mathrm{H} \vetor{v} + \sum_{i}\alpha_{i} \rho_{i} \mathrm{H}_{i} \vetor{v}_{d,i}\right] \cdot \mathrm{d}\vetor{a}
&= - \int\limits_{\partial V} \vetor{\dot{q}''}  \cdot \mathrm{d}\vetor{a} + \int\limits_{\partial V} \tens{T}\cdot\vetor{v}\,\mathrm{d}\vetor{a} \\
&\quad + \int\limits_{V}\vetor{f}_{b}\cdot\vetor{v}\,\mathrm{d}V
+ \int\limits_{V} \mathrm{S}_{\mathrm{E}}\,\mathrm{d}V
\end{split}
\end{equation}

\noindent where
\begin{itemize}[topsep=0pt, labelindent=40pt,labelsep=5pt, parsep = 2pt, listparindent=2em]
\item $\mathrm{E}$ is the total energy,
\item $\mathrm{H}$ is the total enthalpy,
\item $\vetor{\dot{q}''}$ is the heat flux vector,
\item $\tens{T}$ is the viscous stress tensor,
\item $\mathrm{S}_{\mathrm{E}}$ is the user defined energy source term.
\end{itemize}
\subsection{Modelling Turbulence}
Flows are generally turbulent. Turbulence is chracterised by irregular fluid motion dominated by
 eddies and vortices across a wide range of spatial and temporal scales. The largest eddies
 extract kinetic energy from the mean flow, which is then, through a cascade, transferred
 to the eddies at
 the smallest scale where the energy is dissipated due to viscosity, raising the 
 internal energy of the fluid~\cite{gilmanDevelopmentGeneralPurpose2014, davidsonIntroductionTurbulenceModels}.

Although turbulence can be modelled in various ways to various levels of resolutions,
 turbulence in this study has been modelled using the Shear Stress Transport (SST) k-w model.
 The SST k-w model is a two-equation eddy-viscosity turbulence model
 based on the Reynolds-Averaged Navier-Stokes (RANS) equations.
 The model combines the advantages of the $k$-$\omega$ model near walls—which
 accurately resolves boundary layer behaviour—and the $k$-$\varepsilon$ model in
 the free stream, enhancing robustness and accuracy in separated flows~\cite{wilcoxTurbulenceModelingCFD2006}.
 The SST approach uses a blending function to smoothly transition between these models,
 which is especially effective in flows with strong adverse pressure gradients and separation,
 such as multiphase cavity flows~\citelater{Menter1994}.

RANS models result from decomposing flow variables into
 mean and fluctuating components,$\phi = \bar{\phi} + \phi'$, applying Reynolds averaging,
 which filters out turbulent fluctuations below the
 averaging scale~\cite{besnardTurbulenceMultiphaseFlow1988a, hirschNumericalComputationInternal2007}.
 This operation yields additional Reynolds stress terms representing
 momentum transport by turbulence, which require closure.
 In order to draw closure relations, eddy-viscosity models of turbulence assume that:
 \begin{itemize}
    \item the turbulent scalar flux vector aligns with the mean scalar gradient.
    \item the Reynolds stress tensor is proportional to the mean velocity strain rate through an eddy viscosity.
 \end{itemize}
 While these assumptions simplify modelling, they tend to be, more often than not, invalid
 as they do not fully capture turbulence physics, especiialy in anisotropic 
 or non-equilibrium flow~\cite{popeTurbulentFlows2015}. This limitation can be significant
 in regions with complex strain rates or strong secondary flows.

The implementation of the model in Star-CCM+ can be further controlled to
 adjust the realisability. Realisability is the postive semi-definiteness of the 
 reynolds stress tensor. The specification of turbulence is done using
 the turbulence intensity and the length scale. The inherent anisotropy 
 in turbulence is accounted for with quadratic and cubic constitutive options
 in the modelling of reynolds stresses. The parameters in the model
 are left on their default values as they are empirically modelled.

\subsection{Evaporation}
The evaporation process is distinguished between occuring at saturation temperatures or not. 
 Evaporation at saturation temperatures, known as vaporisation, is a much faster 
 process in comparison to evaporation at normal~\cite{wickertSimulation}. 
 In case of vaporisation along a heated wall, most CFD models of
 vaporisation treat the energy transfer by partitioning
 the heat flux at the wall into the latent heat of the bubble vaporisation process and the 
 sensible heat transferred to the liquid through conduction, convection, and quencing processes.
 This is determined by the area fractions covered by the liquid and vapour. Correlations between the
 heat flux and wall superheat~\COMMENT{what is superheat} have been made using empirical
 results or based on modelling of sub-processes.

The sub-processes that are primarily taken into account in such determinations are:
\begin{itemize}
    \item number density of nucleation sites
    \item diameter of bubble at time of detachment from the surface
    \item frequency of bubble detachment
\end{itemize}
Nucleation affects the number and distribution of bubbles in the fluid. Walls, particularly so in case they're heated, are more suspectible to be nucleation sites for bubble growth
 as the activation energy for bubble growth is much lower. This nucleation can be modelled through the domain as homogenous or heterogenous - where nucleation sites are modelled using a 
 presumed PDF function or be determined by empirical correlations.~\cite{liaoComputationalModellingFlash2017}

After a stable bubble nucleus has been formed, the growth of the bubble becomes pertinent in the evaluation
 of the diameter at departure and the interphase energy and mass transfer. In the case of an ideal
 spherical bubble undergoing expansion in a uniformly superheated liquid, the rate of growth
 is determined by the surface tension, liquid inertia, and the differnce in pressures inside
 and outside the bubble. Surface tension dominates and restricts the initial bubble growth.
 This initial rapid expansion of the bubble results in the formation of a thin liquid layer separating
 the vapour from the wall surface. This microlayer can be a significant contributor to the 
 bubble growth and heat transfer process.
 As the radius increases and the liquid layer around the bubble cools due to evaporation,
 a decrease in vapour pressure in seen. The dynamics are controlled by thermal diffusion
 and inertia. Inertial effects, however, become less relevant as the bubble growth rate decreases
 and the bubble radius grows. The interphase energy exchange determines further bubble growth.~\cite{sherFlashboilingAtomization2008}

\subsubsection{Effects of low pressure and gravity conditions}
The absence of significant gravitaional forces alters the dynamics of vaporisation - affecting natural convection, buoyancy
 driven phenomenon, and bubble dynamics. Under terrestrial gravity, convection currents develop under thermal
 gradients, thus bringing cooler liquid to the heated surface as the vapour bubble rises due
 to buyoancy. This density driven natural convection is affected in low gravity conditions, maintaining
 the thermal gradient. Bulk fluid movement is affected and heat and mass transfer processes shift
 from being convection dominated to being diffusion dominated.~\cite{cochranEffectsFluidProperties, leeExperimentalComputationalInvestigation2022}

Under low gravity conditions, vapour bubbles tend to remain near the heat source, growing
 signifcantly larger and coalescing with smaller bubbles. Prolonged bubble presence on the
 surface results in drying out of the surface and unsteady rise in surface temperatures.\cite{qiuSingleBubbleDynamicsPool2002}~\COMMENT{Mention some numbers}
 In conditions where the dominance of surface tension forces becomes more pronounced,
 the surfae tension gradients can induce fluid circulation. This is known as Marangoni
 convection and is observed to affect bubble growth and heat transfer.~\cite{xuNumericalStudyMarangoni2014}

Reduced system pressure impacts the overall intensity of heat transfer. The onset of nucleate boiling also
 requires higher superheat values in low pressure conditions.
 Decrease in pressure leads to a lower number of active nucleation sites and increases
 microlayer thickness found at the bubble-wall interface. The local heat flux density from
 microlayer evaporation is impacted and the microlayer evaporation area and depletion time is
 increased. The system is also susceptible to the formation of a continuous vapour film at the 
 surface due to coalescing bubbles unable to escape the surface, reducing the overall heat transfer.~\cite{serdyukovInfluencePressureLocal2023}
 
The affected heat transfer and bubble growth phenomenon further result in lower critical heat fluxes
 and heat transfer efficiency.~\cite{zhangExperimentalTheoreticalStudy2002, PDFMethodAssessing}

\subsubsection{Vaporisation in a superheated liquid}
\COMMENT{------------UNDER CONSTRUCTION--------------------}
Flash evaporation is the intense vaporisation triggered through depressurisation of superheated fluids.

\subsection{Multiple Flow Regimes and Interface Capturing}
Star-CCM+ in its implementation of the eulerian-multiphase framework utilises a
 dynamic flow regime topology based on the local volume fractions
 of the interacting phases. Based on the increasing volume fraction of the
 secondary phase in a computational cell,
 the topology transitions through three principal regimes:
 \begin{itemize}
    \item dispersed in the primary phase
    \item having a large scale interface with the primary phase
    \item continuous while having the primary phase dispersed
 \end{itemize}
 For each cell, Star-CCM+ uses regime-specific closure models for the interfacial~\COMMENT{OR RATHER INTERPHASE MOMENTUM AND ENERGY TRANSFER}
 forces.
 These models are blended using weighting schemes 
 according to the calculated proportion of each regime in the cell.~\cite{COLOMBO2019968}
 ~\citelater{Gada et al. (2015): Large Scale Interface Model for Simulating Multiphase Flows (STAR-CCM+ Implementation)}

\COMMENT{EDIT}While there are multiple numerical methods to account for volume fraction transport, the 
 High-Resolution Interface Capturing scheme is specifically implemented to handle 
 sharp phase interfaces. Free-surface flows with sharp 
 interfaces present a discontinuity
 in the solution domain, causing discretisation errors. This can lead to
 spurious or parasitic currents.~\cite{MIRJALILI2019221}
 Although often negligible, these currents can be accommodated for through 
 the incorporation of an artificial viscosity term.   

At large scale interfaces, the near-interface flow often exhibits
 laminar characteristics with the development of boundary layers. However, in situations
 with significant velocity gradients, turbulent fluctuations can arise in the
 regions adjacent to the interface (especially when the interface is mobile or
 highly distorted). Recognizing that standard RANS turbulence
 models tend to overpredict turbulence intensity near phase boundaries,
 STAR-CCM+ allows the user to apply turbulence damping strategies
 (e.g., vorticity limiters) which suppress excessive eddy viscosity and
 restore physically consistent laminar sublayers adjacent to the interface.

\subsection{Phase Interactions}
\COMMENT{------------UNDER CONSTRUCTION--------------------}
Phase interactions in free-surface flows and flash evaporation dictate
 how mass, momentum and energy are exchanged between contiguous liquid and
 gaseous domains. For both separated (e.g., stratified) and dispersed (e.g., bubbly) flow regimes, 
major phase interactions include:
\begin{itemize}
    \item Drag force:
     Governs momentum exchange due to the relative motion between phases;
     typically decomposed using empirical drag laws (e.g., Schiller-Naumann for
     bubbles/droplets).
    \item Lift force:
     Acts perpendicular to the relative flow, important in shear
flows with small dispersed phases; affects lateral migration of bubbles/droplets—uses
models like Tomiyama, Saffman, or others depending on flow characteristics.
    \item Virtual mass force (added mass)
     Represents resistance to acceleration caused
by the acceleration of surrounding fluid (significant in cases with rapid
acceleration or deceleration of bubbles).
    \item Wall lubrication force
     Repels dispersed phase from solid walls, preventing unrealistic
accumulation at boundaries (important in vertical/horizontal pipe flows).
    \item Turbulent dispersion force
     Spreads the dispersed phase due to turbulence
in the continuous phase, modeled typically as a diffusion term proportional
to turbulent kinetic energy gradient.
 The Favre-Averaged Drag (FAD) model
 approximates turbulent dispersion as a function of liquid
 eddy viscosity and bubble slip. This mitigates excessive bubble
 clustering below the free surface.~\cite{liaoFlashingEvaporationDifferent2013}
    \item Surface tension
     Maintains interface shape, forms basis for capillary effects
and controls droplet/bubble formation and coalescence/breakup.
    \item Interphase mass transfer (evaporation/condensation)
     Represents phase change, modeled as interfacial flux determined
by local thermodynamic properties (temperature, concentration, pressure—often
using kinetic models for flash evaporation).
    \item Interphase energy transfer
     Encompasses latent heat and sensible heat exchange at the interface 
and throughout the phases (essential for accurate phase change modeling).
    \item Nucleation (wall and bulk)
    \item Interfacial area changes due to bubble growth, coalescence, breakup
    \item Local pressure and temperature gradients driving simultaneous mass and energy transfer
\end{itemize}
\subsection{Numerical Methods}
\COMMENT{------------UNDER CONSTRUCTION--------------------}
\begin{itemize}
    \item gradients
    \item solvers - implicit unsteady, segregated flow (SIMPLE)
    \item Relevant non-dim numbers? (reynolds, froude, jakobs, nusselt, bond, what not)
\end{itemize}
\subsection{Workflow?}
\COMMENT{------------UNDER CONSTRUCTION--------------------}